\documentclass[10pt,a4paper]{amsart}
\usepackage{amsmath,amssymb,amsthm,graphicx}
% Retained minimal preamble for the Stacks Project chapter more-morphisms.tex (More on Morphisms).
% Verbatim excerpt of the chapter preamble: Xy-pic setup (native preamble lines 14-18) and the native
% theorem declarations with the mathematical macros (native preamble lines 154-217). Copied mechanically;
% nothing added, removed or redefined. The wrapper class and the amsmath/amssymb/amsthm/graphicx base
% packages are supplied by the assembler, so the class/externaldocument part of the native preamble is omitted.
% For commutative diagrams we use Xy-pic
\usepackage[all]{xy}

% We use 2cell for 2-commutative diagrams.
\xyoption{2cell}
% Theorem environments.
%
\theoremstyle{plain}
\newtheorem{theorem}[subsection]{Theorem}
\newtheorem{proposition}[subsection]{Proposition}
\newtheorem{lemma}[subsection]{Lemma}

\theoremstyle{definition}
\newtheorem{definition}[subsection]{Definition}
\newtheorem{example}[subsection]{Example}
\newtheorem{exercise}[subsection]{Exercise}
\newtheorem{situation}[subsection]{Situation}

\theoremstyle{remark}
\newtheorem{remark}[subsection]{Remark}
\newtheorem{remarks}[subsection]{Remarks}

\numberwithin{equation}{subsection}

% Macros
%
\def\lim{\mathop{\mathrm{lim}}\nolimits}
\def\colim{\mathop{\mathrm{colim}}\nolimits}
\def\Spec{\mathop{\mathrm{Spec}}}
\def\Hom{\mathop{\mathrm{Hom}}\nolimits}
\def\Ext{\mathop{\mathrm{Ext}}\nolimits}
\def\SheafHom{\mathop{\mathcal{H}\!\mathit{om}}\nolimits}
\def\SheafExt{\mathop{\mathcal{E}\!\mathit{xt}}\nolimits}
\def\Sch{\mathit{Sch}}
\def\Mor{\mathop{\mathrm{Mor}}\nolimits}
\def\Ob{\mathop{\mathrm{Ob}}\nolimits}
\def\Sh{\mathop{\mathit{Sh}}\nolimits}
\def\NL{\mathop{N\!L}\nolimits}
\def\CH{\mathop{\mathrm{CH}}\nolimits}
\def\proetale{{pro\text{-}\acute{e}tale}}
\def\etale{{\acute{e}tale}}
\def\QCoh{\mathit{QCoh}}
\def\Ker{\mathop{\mathrm{Ker}}}
\def\Im{\mathop{\mathrm{Im}}}
\def\Coker{\mathop{\mathrm{Coker}}}
\def\Coim{\mathop{\mathrm{Coim}}}

% Boxtimes
%
\DeclareMathSymbol{\boxtimes}{\mathbin}{AMSa}{"02}

%
% Macros for moduli stacks/spaces
%
\def\QCohstack{\mathcal{QC}\!\mathit{oh}}
\def\Cohstack{\mathcal{C}\!\mathit{oh}}
\def\Spacesstack{\mathcal{S}\!\mathit{paces}}
\def\Quotfunctor{\mathrm{Quot}}
\def\Hilbfunctor{\mathrm{Hilb}}
\def\Curvesstack{\mathcal{C}\!\mathit{urves}}
\def\Polarizedstack{\mathcal{P}\!\mathit{olarized}}
\def\Complexesstack{\mathcal{C}\!\mathit{omplexes}}
% \Pic is the operator that assigns to X its picard group, usage \Pic(X)
% \Picardstack_{X/B} denotes the Picard stack of X over B
% \Picardfunctor_{X/B} denotes the Picard functor of X over B
\def\Pic{\mathop{\mathrm{Pic}}\nolimits}
\def\Picardstack{\mathcal{P}\!\mathit{ic}}
\def\Picardfunctor{\mathrm{Pic}}
\def\Deformationcategory{\mathcal{D}\!\mathit{ef}}
\newenvironment{reference}{}{}
\newenvironment{slogan}{}{}
\newenvironment{history}{}{}
\usepackage{cleveref}
\crefname{originallemma}{original lemma}{original lemmas}
\Crefname{originallemma}{Original lemma}{Original lemmas}
\makeatletter
\newlabel{algebra-lemma-quasi-finite-open-integral-closure}{{00QB}{source}{}{sourcepointer-1}{}}
\newlabel{algebra-lemma-quasi-finite-open-integral-closure@cref}{{[originallemma][0][]00QB}{[1][0][]source}{}{}{}}
\makeatother
\crefname{originallemma}{original lemma}{original lemmas}
\Crefname{originallemma}{Original lemma}{Original lemmas}
\makeatletter
\newlabel{lemma-Noetherian-approximation-dimension-d}{{05FJ}{source}{}{sourcepointer-2}{}}
\newlabel{lemma-Noetherian-approximation-dimension-d@cref}{{[originallemma][0][]05FJ}{[1][0][]source}{}{}{}}
\makeatother
\crefname{originallemma}{original lemma}{original lemmas}
\Crefname{originallemma}{Original lemma}{Original lemmas}
\makeatletter
\newlabel{lemma-open-and-closed-in-finite}{{0F3M}{source}{}{sourcepointer-3}{}}
\newlabel{lemma-open-and-closed-in-finite@cref}{{[originallemma][0][]0F3M}{[1][0][]source}{}{}{}}
\makeatother
\crefname{originallemma}{original lemma}{original lemmas}
\Crefname{originallemma}{Original lemma}{Original lemmas}
\makeatletter
\newlabel{morphisms-lemma-locally-quasi-finite-rel-dimension-0}{{0397}{source}{}{sourcepointer-4}{}}
\newlabel{morphisms-lemma-locally-quasi-finite-rel-dimension-0@cref}{{[originallemma][0][]0397}{[1][0][]source}{}{}{}}
\makeatother
\title{Universally parametrizing the open-and-closed subschemes of a quasi-finite map that are finite over the base}
\author{The Stacks Project Authors}
\date{stacks-project:more-morphisms.tex:local-snapshot-20261011; GNU FDL 1.2 or later}
\begin{document}
\maketitle
\noindent\textit{Editorial note.} Statement span is lines 22834-22845 of sources/stacks-more-morphisms/pinned-original-source.tex; proof span is lines 22847-22888, the complete author proof (40 body lines).\par\medskip
\noindent\textit{Editorial note.} Navigation stubs: 4 cited-but-not-included labels are supplied as non-load-bearing pointers with their pinned tags where the snapshot tag index has them.\par\medskip
\noindent\textit{Editorial note.} extra\_packages is empty on purpose: the retained preamble already loads xy with options and a second declaration clashes in the assembler wrapper.\par\medskip
\noindent\textit{Editorial note.} Text is retained verbatim, including the author's own typos and placeholders.\par\medskip
\noindent\textit{Editorial note.} The quasi-finite companion of the chapter's finite-case moduli-of-open-and-closed-subschemes lemma; the author's automatic-closedness sentence for open immersions between schemes finite over Y' is the only inline-implicit step and is a one-line standard fact.\par\medskip
\noindent\textit{Allowed standard prerequisites (editorial).}\begin{itemize}
\item Explicit prerequisite boundary. Morphisms, lemma locally-quasi-finite-rel-dimension-0; the chapter's Noetherian-approximation-dimension-d and open-and-closed-in-finite lemmas (both previously collected); Algebra, lemma quasi-finite-open-integral-closure (quasi-finite morphisms are opens of finite ones after integral closure); standard idempotent/open-and-closed dictionary.
\end{itemize}\medskip
\ifdefined\hypertarget\hypertarget{sourcepointer-1}{}\fi
\noindent\textit{Source reference.} Original lemma 00QB: sources/stacks-more-morphisms/pinned-original-source.tex, source lines 22888--22888. Cited by the retained statement or proof but not part of the collected spans; supplied only as a navigation pointer. The argument does not depend on its content.\par\medskip
\ifdefined\hypertarget\hypertarget{sourcepointer-2}{}\fi
\noindent\textit{Source reference.} Original lemma 05FJ: sources/stacks-more-morphisms/pinned-original-source.tex, source lines 9950--9957. Cited by the retained statement or proof but not part of the collected spans; supplied only as a navigation pointer. The argument does not depend on its content.\par\medskip
\ifdefined\hypertarget\hypertarget{sourcepointer-3}{}\fi
\noindent\textit{Source reference.} Original lemma 0F3M: sources/stacks-more-morphisms/pinned-original-source.tex, source lines 22791--22797. Cited by the retained statement or proof but not part of the collected spans; supplied only as a navigation pointer. The argument does not depend on its content.\par\medskip
\ifdefined\hypertarget\hypertarget{sourcepointer-4}{}\fi
\noindent\textit{Source reference.} Original lemma 0397: sources/stacks-more-morphisms/pinned-original-source.tex, source lines 22888--22888. Cited by the retained statement or proof but not part of the collected spans; supplied only as a navigation pointer. The argument does not depend on its content.\par\medskip
\begin{lemma}
\label{lemma-open-and-closed-in-quasi-finite}
Let $X \to Y$ be a morphism of affine schemes which is quasi-finite and
of finite presentation. There exists a morphism $Y_{univ} \to Y$
of finite presentation and an open subscheme
$U_{univ} \subset Y_{univ} \times_Y X$ such that
$U_{univ} \to Y_{univ}$ is finite with the following property:
given any morphism $Y' \to Y$ of affine schemes
and an open subscheme $U' \subset Y' \times_Y X$
such that $U' \to Y'$ is finite, there exists a morphism
$Y' \to Y_{univ}$ such that the inverse image of $U_{univ}$ is $U'$.
\end{lemma}

\begin{proof}
Recall that a finite type morphism is quasi-finite if and only
if it has relative dimension $0$, see
Morphisms, Lemma \ref{morphisms-lemma-locally-quasi-finite-rel-dimension-0}.
By Lemma \ref{lemma-Noetherian-approximation-dimension-d}
applied with $d = 0$ we reduce to the case where
$X$ and $Y$ are Noetherian. We may choose an open immersion
$X \to X'$ such that $X' \to Y$ is finite, see
Algebra, Lemma \ref{algebra-lemma-quasi-finite-open-integral-closure}.
Note that if we have $Y' \to Y$ and $U'$ as in (2), then
$$
U' \to Y' \times_Y X \to Y' \times_Y X'
$$
is open immersion between schemes finite over $Y'$ and hence
is closed as well. We conclude that $U'$ corresponds to an
idempotent in
$$
\Gamma(Y', \mathcal{O}_{Y'})
\otimes_{\Gamma(Y, \mathcal{O}_Y)}
\Gamma(X', \mathcal{O}_{X'})
$$
whose corresponding open and closed subset is contained in
the open $Y' \times_Y X$. Let $Y'_{univ} \to Y$ and idempotent
$$
e'_{univ} \in
\Gamma(Y_{univ}, \mathcal{O}_{Y_{univ}})
\otimes_{\Gamma(Y, \mathcal{O}_Y)}
\Gamma(X', \mathcal{O}_{X'})
$$
be the pair constructed in Lemma \ref{lemma-open-and-closed-in-finite}
for the ring map $\Gamma(Y, \mathcal{O}_Y) \to \Gamma(X', \mathcal{O}_{X'})$
(here we use that $Y$ is Noetherian to see that $X'$ is of finite presentation
over $Y$). Let $U'_{univ} \subset Y'_{univ} \times_Y X'$ be the corresponding
open and closed subscheme. Then we see that
$$
U'_{univ} \setminus Y'_{univ} \times_Y X
$$
is a closed subset of $U'_{univ}$ and hence has closed image
$T \subset Y'_{univ}$. If we set $Y_{univ} = Y'_{univ} \setminus T$
and $U_{univ}$ the restriction of $U'_{univ}$ to
$Y_{univ} \times_Y X$, then we see that the lemma is true.
\end{proof}

\end{document}
